% Part: sets-functions-relations
% Chapter: functions
% Section: functions-relations

\documentclass[../../../include/open-logic-section]{subfiles}


\begin{document}

\olfileid{sfr}{fun}{rel}

\olsection{Fungsi minangka Relasi}

\begin{explain}
Fungsi kang metakake !!{element}s saka~$A$ menyang !!{element}s saka~$B$
cetha nemtokake relasi antarane $A$ lan~$B$, yaiku relasi
kang dumadi antarane $x$ lan $y$ yen lan mung yen $f(x) = y$. Malah,
yen kita kepengin nyuda cacah jinis bahan dhasar matematika, upamane,
kita bisa \emph{nganggep padha} fungsi~$f$ lan relasi iki,
yaiku himpunan pasangan. Iki banjur nuwuhake pitakon:
relasi endi wae kang nemtokake fungsi kanthi cara iki?
\end{explain}

\begin{defn}[Graf fungsi] Upamakna $f\colon A \to B$ sawijining fungsi.
\emph{Graf} saka~$f$ yaiku relasi $R_f \subseteq A \times B$
kang ditetepake kanthi
\[
R_f = \Setabs{\tuple{x,y}}{f(x) = y}.
\]
\end{defn}

\begin{explain}
Adhedhasar ekstensionalitas, graf sawijining fungsi ditemtokake kanthi tunggal.
Kajaba iku, ekstensionalitas tumrap himpunan langsung mbenerake
prinsip ekstensionalitas fungsi kang sadurunge dianggep lumaku:
yen $f$ lan~$g$ duwe domain lan kodomain kang padha,
fungsi loro iku padha yen kabeh nilaie padha.

Semono uga, yen sawijining relasi iku "fungsional", relasi iku graf sawijining fungsi.
\end{explain}

\begin{prop}\ollabel{prop:graph-function}
Upamakna $R \subseteq A \times B$ nduweni sipat:
\begin{enumerate}
\item Yen $Rxy$ lan $Rxz$, mula $y = z$; lan
\item kanggo saben $x \in A$ ana $y \in B$ saengga $\tuple{x,
y} \in R$.
\end{enumerate}
Banjur $R$ iku graf fungsi $f\colon A \to B$ kang ditetepake kanthi
$f(x) = y$ yen lan mung yen $Rxy$.
\end{prop}

\begin{proof}
Upamakna ana $y$ saengga $Rxy$. Yen ana $z \neq
y$ liyane saengga $Rxz$, sarat tumrap~$R$ bakal dilanggar. Mula, yen
ana $y$ saengga $Rxy$, $y$ iki tunggal, saengga $f$
ditetepake kanthi becik. Cetha yen $R_f = R$.
\end{proof}

\begin{explain}
Saben fungsi $f\colon A \to B$ duwe graf, yaiku relasi kang kaemot ing $A
\times B$ lan ditetepake dening $f(x) = y$. Cathetan koreksi sumber
OLFUN-004: iki relasi antarane A lan B, dudu relasi ing himpunan
produk A lan B miturut definisi relasi biner ing buku iki. Kosok baline, saben relasi~$R
\subseteq A \times B$ kang nduweni sipat ing
\olref{prop:graph-function} iku graf sawijining fungsi~$f \colon A \to
B$. Amarga sesambungan kang raket antarane fungsi lan grafe iki,
kita bisa nganggep sawijining fungsi minangka grafe. Kanthi tembung
liya, fungsi bisa dianggep padha karo relasi tartamtu, yaiku
himpunan tupel tartamtu. \oliflabeldef{sfr:rel:ref:sec}{Nanging, gatekna
yen maksud "nganggep padha" iki kaya ing
\olref[sfr][rel][ref]{sec}: iki dudu pratelan ngenani metafisika
fungsi, nanging pangamatan yen \emph{nganggep}
fungsi minangka himpunan tartamtu iku migunani. Salah siji sebabe
yaiku saiki k}{K}ita bisa nindakake operasi ing
fungsi kang padha karo operasi ing relasi (delengen
\olref[sfr][rel][ops]{sec}). Mligine:
\end{explain}

\begin{defn}\ollabel{defn:funimage}
Upamakna $f \colon A \to B$ sawijining fungsi kanthi $C\subseteq A$.

\emph{Watesan} saka~$f$ ing~$C$ yaiku
fungsi~$\funrestrictionto{f}{C}\colon C \to B$ kang ditetepake kanthi
$(\funrestrictionto{f}{C})(x) = f(x)$ kanggo kabeh $x \in C$. Kanthi tembung
liya, $\funrestrictionto{f}{C} = \Setabs{\tuple{x, y} \in R_f}{x \in
C}$.

\emph{Aplikasi} saka~$f$ tumrap~$C$ yaiku $\funimage{f}{C} =
\Setabs{f(x)}{x \in C}$. Iki uga diarani \emph{citra} saka~$C$
dening~$f$.
\end{defn}

\begin{explain}
Saka definisi iki bisa didudut yen $\ran{f} =
\funimage{f}{\dom{f}}$, kanggo saben fungsi~$f$.
\oliflabeldef{sfr:rel:ops:sec}{Pangerten iki analog karo
definisi ing \olref[sfr][rel][ops]{sec} lan anggone kita
nganggep fungsi padha karo relasi, nanging ora persis padha:
watesan fungsi ing kene mung nyaring koordinat masukan, dene
watesan relasi ing kana nyaring loro koordinat. Iki koreksi
andharan sumber OLFUN-005. Rong operasi liyane,
yaiku invers lan produk relatif, mbutuhake katrangan luwih jangkep.
Katrangan iku bakal diwenehake ing \olref[inv]{sec} lan
\olref[cmp]{sec}.}{}
\end{explain}

\end{document}
